\chapter{Estado del arte}
\label{cap:estado-arte}

En este capítulo se revisarán los principales métodos de segmentación e identificación automática de dientes en IOS 3D. Para ello, se analizarán las diferentes estrategias según la representación de los datos, los procedimientos de numeración y la utilización de la forma de la arcada dental, así como los resultados del reto 3DTeethSeg. Esta revisión permitirá identificar las limitaciones que permanecen abiertas en la literatura y que motivan la necesidad de abordar el problema planteado en el Capítulo~\ref{cap:introduccion} mediante enfoques alternativos.

\section{Segmentación de dientes en IOS 3D}
\label{sec:ea-segmentacion}

En esta sección se clasificarán los métodos de segmentación dental según la representación de los datos de entrada, distinguiendo cuatro familias: imágenes bidimensionales, vóxeles, nubes de puntos y mallas o grafos. Se describirán sus principales características mediante métodos representativos, como FiboSeg, TeethSeg y CGIP \cite{ben-hamadou3DTeethSeg223DTeeth2023}, TSegNet \cite{cuiTSegNetEfficientAccurate2021} y MeshSegNet \cite{lianDeepMultiScaleMesh2020}. Asimismo, se analizarán las ventajas y limitaciones de cada representación, considerando aspectos como la pérdida de información en superficies ocultas, el coste computacional, la resolución espacial y la conservación de la conectividad. Dado que las arquitecturas de aprendizaje profundo ya se han presentado en la Sección~\ref{sec:aprendizaje-3d}, el análisis se centrará en su aplicación específica al problema de la segmentación dental.

\section{Identificación de los dientes}
\label{sec:ea-identificacion}

En esta sección se analizarán las estrategias utilizadas para asignar las etiquetas del sistema FDI a los dientes previamente segmentados. En primer lugar, se describirá el procedimiento de TSegNet \cite{cuiTSegNetEfficientAccurate2021}, que realiza la identificación de cada pieza de forma independiente. En consecuencia, se señalará que este enfoque no garantiza la unicidad de las etiquetas asignadas ni representa explícitamente las posiciones vacantes. En cambio, se estudiará el procedimiento de postprocesado de IGIP \cite{ben-hamadou3DTeethSeg223DTeeth2023}, que utiliza la ordenación de los dientes a lo largo de la curva de la arcada para corregir etiquetas duplicadas y errores de secuencia, aunque no contempla el tratamiento de las posiciones vacantes.

\section{Uso de la forma de la arcada dental}
\label{sec:ea-arcada}

En esta sección se analizará cómo los métodos existentes incorporan la disposición espacial de los dientes como conocimiento previo para mejorar su detección e identificación. En primer lugar, se describirá DArch \cite{qiuDArchDentalArch2022}, que estima la curva de la arcada mediante una aproximación con curvas de Bézier y un refinamiento posterior con una red convolucional de grafos (GCN), utilizándola para orientar el muestreo de puntos y facilitar la detección de centroides, sin realizar la numeración dental. A diferencia de DArch, se estudiará SRF-Net \cite{maSRFNetSpatialRelationship2020}, que incorpora la posición relativa de los dientes y las relaciones de similitud con sus vecinos para mejorar la clasificación, además de utilizar una plantilla para corregir errores de etiquetado y números duplicados. No obstante, se señalará que SRF-Net depende de una distribución dental relativamente regular, tal como indica el equipo IGIP \cite{ben-hamadou3DTeethSeg223DTeeth2023}. Asimismo, se explicará que, ante ausencias que generan ambigüedad, SRF-Net \cite{maSRFNetSpatialRelationship2020} adopta un criterio de clasificación basado en el tipo de diente, sin exigir la identificación exacta de su posición FDI.

\section{El reto 3DTeethSeg}
\label{sec:ea-reto}

En esta sección se presentará el reto 3DTeethSeg \cite{ben-hamadou3DTeethSeg223DTeeth2023}, destacando sus principales aportaciones: el conjunto de datos Teeth3DS, descrito en la Sección~\ref{sec:datos}, las métricas de evaluación definidas en la Sección~\ref{sec:metricas} y un protocolo de evaluación basado en un conjunto de prueba oculto. A continuación, se compararán los resultados de los métodos participantes, prestando especial atención a CGIP, ganador de la clasificación general, e IGIP, que obtuvo el mejor resultado en identificación dental (TIR). Esta comparación permitirá poner de manifiesto que un mayor rendimiento en segmentación no implica necesariamente una mejor identificación, como también refleja el caso de FiboSeg, que incorpora ausencias dentales simuladas durante el entrenamiento y que, además, pasa del cuarto puesto en TSA al segundo en TIR. Asimismo, se recogerán las limitaciones reconocidas por los organizadores, entre ellas la necesidad de incorporar casos con dientes ausentes o dañados y situaciones de etiquetado ambiguo, así como de ampliar la evaluación a la suavidad de los límites entre dientes y encía y al coste computacional de los métodos.

\section{Síntesis}
\label{sec:ea-sintesis}

En esta sección se sintetizarán las principales conclusiones extraídas de la revisión de los métodos existentes. Se pondrá de relieve que, aunque las técnicas actuales alcanzan buenos resultados de segmentación, las estrategias de numeración estudiadas, como TSegNet, IGIP y SRF-Net, presentan limitaciones para garantizar la coherencia global de las etiquetas cuando existen ausencias dentales  (Secciones~\ref{sec:ea-identificacion} y~\ref{sec:ea-arcada}). Asimismo, se señalará que las métricas de evaluación habituales, descritas en la Sección~\ref{sec:metricas}, promedian los resultados sobre el conjunto de dientes y no permiten identificar explícitamente los errores de coherencia asociados a las ausencias. En consecuencia, se planteará el interés de integrar información geométrica y posicional en estrategias de asignación global que permitan preservar la coherencia de la numeración dental.